\documentclass[12pt, ukrainian]{article}
\usepackage[a4paper]{geometry}
\setlength\parindent{0mm}
\setlength\parskip{4mm}
\geometry{top=1.5cm,bottom=1.5cm,left=1.3cm,right=1.5cm}
\pagestyle{empty}
\usepackage[T2A]{fontenc}
\usepackage[utf8]{inputenc}
\usepackage[ukrainian]{babel}
\usepackage{graphicx}
\usepackage{caption}
\usepackage{caption}
\DeclareCaptionFormat{nocaption}{}
\captionsetup{format=nocaption,aboveskip=0pt,belowskip=0pt}
\usepackage[Export]{adjustbox}
\adjustboxset{max size={0.9\linewidth}{0.9\paperheight}}
\usepackage{verbatim, listings, clrscode3e, fancyvrb}
\def\code#1{\Verb!#1!}
\begin{document}
\begin {large}

\textbf{01.12.2023 Контрольна робота, варіант  \No { 22 }}\par


{
{\begin {center}\textbf{Завдання 1}\end {center}}\par
По яких днях тижня відбувались практичні з предмету?\par
{\begin {center}\textbf{Завдання 2}\end {center}}\par
Як зовуть (ім'я, по батькові) викладача, що вів практичні заняття?\par
{\begin {center}\textbf{Завдання 3}\end {center}}\par
Виберіть всі правильні твердження, якщо такі тут є.\par
\vspace{2mm}
( 1 ) Деякі задачі про перестановки з забороненими позиціями можуть бути розв’я\-за\-ні за допомогою формули включення та виключення.%good

( 2 ) Якщо дошка складається з кількох частин, що не мають спільних рядків та спільних стовпчиків, то її туровий поліном можна отримати як добуток поліномів її частин.%good

( 3 ) За означенням туровий поліном є експоненційною твірною функцією.%bad

( 4 ) Послідовність, що має експоненційну твірну, звичайної твірної може не мати.%good

( 5 ) Турові поліноми --- це ані про тури, ані про поліноми.%bad

( 6 ) Задачу знаходження кількості способів сформувати новорічні подарункові набори з $n$ цукерок, що складаються з цукерок 13 видів, де цукерок кожного вигляду не більше 20, можна розв’язати за допомогою експоненційних твірних.%bad

( 7 ) Турові поліноми можна розглядати не тільки для квадратних дошок.%good

( 8 ) Послідовність, задана рекурентним співвідношенням $a_{n+2}=a_{n+1}+a_{n}$ може не бути послідовністю чисел Фібоначчі.%good

( 9 ) Турові поліноми розглядаються виключно для квадратних дошок.%bad

( 10 ) Кожна числова послідовність має не більше однієї твірної.%good\par
}
\end {large}

\vspace{4mm}
%end
\newpage
\end{document}

